\documentclass[11pt]{article}

\usepackage[utf8]{inputenc}
\usepackage[T1]{fontenc}
\usepackage[spanish,es-tabla]{babel}
\usepackage{lmodern}
\usepackage[a4paper,margin=2cm]{geometry}
\usepackage{graphicx}
\usepackage{xcolor}
\usepackage{colortbl}
\usepackage{array}
\usepackage{enumitem}
\usepackage{tikz}
\usetikzlibrary{backgrounds}
\usepackage{longtable}
\usepackage{booktabs}
\usepackage{hyperref}
\hypersetup{colorlinks=true,linkcolor=black,urlcolor=[HTML]{0B3A5C}}

\graphicspath{{capturas/}}

% Enlaces al prototipo publicado en GitHub Pages.
\newcommand{\urlapp}{https://thankstochatgpt.github.io/LOARD}
\newcommand{\app}[1]{\href{\urlapp#1}{\texttt{#1}}}
\setlength{\parskip}{0.5em}
\setlength{\parindent}{0pt}
\renewcommand{\arraystretch}{1.35}
\setlist{nosep,leftmargin=1.2em}

\definecolor{marca}{HTML}{0B3A5C}
\definecolor{borde}{HTML}{B8C4CE}
\definecolor{barra}{HTML}{E6ECF1}
\definecolor{telefono}{HTML}{1F2933}
\definecolor{suave}{HTML}{F3F6F9}
\definecolor{ok}{HTML}{1E7B45}
\arrayrulecolor{borde}

% ---------- Marcos para capturas (mismos del documento de diseño responsivo) ----------
% Captura de escritorio dentro de una ventana de navegador. #1 ancho, #2 archivo
\newcommand{\escritorio}[2]{%
  \begin{tikzpicture}[baseline=(img.south)]
    \node[inner sep=0] (img) {\includegraphics[width=#1]{#2}};
    \fill[barra] ([yshift=0.36cm]img.north west) rectangle (img.north east);
    \foreach \i/\c in {0/red!55, 1/orange!65, 2/green!50!black!45}
      \fill[\c] ([xshift={0.22cm+\i*0.24cm}, yshift=0.18cm]img.north west) circle (0.06cm);
    \draw[borde, line width=0.8pt, rounded corners=3pt]
      ([yshift=0.36cm]img.north west) rectangle (img.south east);
  \end{tikzpicture}}

% Captura de móvil dentro de un marco de teléfono. #1 alto, #2 archivo
\newcommand{\movil}[2]{%
  \begin{tikzpicture}[baseline=(img.south)]
    \node[inner sep=0] (img) {\includegraphics[height=#1]{#2}};
    \begin{scope}[on background layer]
      \fill[telefono, rounded corners=9pt]
        ([shift={(-0.14cm,0.3cm)}]img.north west) rectangle ([shift={(0.14cm,-0.3cm)}]img.south east);
    \end{scope}
    \fill[telefono!45] ([yshift=0.15cm]img.north) circle (0.045cm);
  \end{tikzpicture}}

% Recorte de pantalla con borde fino. #1 ancho, #2 archivo
\newcommand{\recorte}[2]{%
  \begin{tikzpicture}[baseline=(img.south)]
    \node[inner sep=0] (img) {\includegraphics[width=#1]{#2}};
    \draw[borde, line width=0.6pt, rounded corners=2pt] (img.north west) rectangle (img.south east);
  \end{tikzpicture}}

% Pie corto debajo de cada captura.
\newcommand{\pie}[1]{{\small\color{marca}\textbf{#1}}}

% ---------- Muestras de color ----------
% Texto de ejemplo. #1 color del texto, #2 color del fondo, #3 texto
\newcommand{\muestra}[3]{{\setlength{\fboxsep}{3pt}\fcolorbox[HTML]{B8C4CE}{#2}{\makebox[2.9cm][l]{\textcolor[HTML]{#1}{\small\textbf{#3}}}}}}
% Borde o anillo de ejemplo. #1 color del borde, #2 color del fondo
\newcommand{\muestraborde}[2]{{\setlength{\fboxsep}{4pt}\colorbox[HTML]{#2}{\makebox[2.8cm][c]{\setlength{\fboxrule}{1.8pt}\setlength{\fboxsep}{0pt}\fcolorbox[HTML]{#1}{#2}{\rule{0pt}{8pt}\hspace{1.6cm}}}}}}

% ---------- Resultados ----------
\newcommand{\punto}[1]{\tikz[baseline=-0.6ex]\fill[#1] (0,0) circle (0.1);}
\newcommand{\cumple}{\punto{ok}~\textcolor{ok}{\textbf{Cumple}}}

% Recuadro gris con título.
\newcommand{\recuadro}[2]{%
  \par\smallskip\noindent
  \colorbox{suave}{\begin{minipage}{\dimexpr\linewidth-2\fboxsep\relax}
    \vspace{4pt}\textbf{#1}\par\smallskip
    #2\vspace{4pt}
  \end{minipage}}\par}

% Tecla dibujada.
\newcommand{\tecla}[1]{\tikz[baseline=(k.base)]\node[draw=borde, fill=suave, rounded corners=2pt, inner sep=2pt, font=\small\sffamily] (k) {#1};}

\title{Auditoría de accesibilidad (WCAG 2.1 AA)\\[0.3em]
\large Loard - Hito 2 - Módulo III: QC de usabilidad y accesibilidad}
\author{}
\date{14 de noviembre de 2026 - actualizada el 17 de noviembre de 2026}

\begin{document}
\maketitle
\thispagestyle{empty}

\section{Introducción}

Se revisó si el prototipo de Loard se puede \textbf{leer bien} y \textbf{usar sin ratón}. La revisión se hizo sobre el prototipo en funcionamiento, siguiendo las pautas de accesibilidad web \textbf{WCAG 2.1, nivel AA}, que es el nivel que se pide normalmente a sitios y aplicaciones. El prototipo revisado está publicado en \href{\urlapp/}{\nolinkurl{https://thankstochatgpt.github.io/LOARD/}}.



Loard no se usa solo en una oficina. Buena parte del trabajo se hace en la obra, y eso hace que estos tres puntos sean necesarios:

\begin{itemize}
  \item \textbf{Mucho contraste por la luz del sol.} En la obra la pantalla del celular o del portátil se usa a pleno sol. La luz solar opaca la pantalla y los colores parecidos entre sí se pierden. Con un contraste alto entre la letra y el fondo, la información se sigue leyendo aunque haya mucha luz.
  \item \textbf{Letras de buen tamaño para leer bien.} Cantidades, fechas, horas y estados de los trabajadores tienen que leerse sin esfuerzo y sin errores. Además, los entornos de obra no facilitan el uso de gafas: el casco, la careta, las gafas de seguridad, el polvo y el sudor las hacen incómodas. Y la soldadura, por la luz intensa del arco, suele afectar la visión con el tiempo.
  \item \textbf{Navegación por teclado para trabajar más rápido.} El residente y la coordinación SST usan el programa durante muchas horas para registrar cortes, personal y cursos. Pasar de un campo a otro con \tecla{Tab}, marcar con \tecla{Espacio} y guardar con \tecla{Enter} es más rápido que mover el ratón para cada campo, y agiliza el trabajo diario.
\end{itemize}

\subsection*{Resultados en breve}

\begin{center}
\begin{tabular}{|>{\bfseries\raggedright\arraybackslash}m{3.4cm}|>{\raggedright\arraybackslash}m{2.4cm}|>{\raggedright\arraybackslash}m{9.8cm}|}
\hline
\rowcolor{barra}
\textbf{Aspecto} & \textbf{Resultado} & \textbf{En pocas palabras} \\ \hline
Contraste de color & \cumple & Se midieron 1.056 textos en 9 pantallas y en los dos temas (claro y oscuro). Ninguno quedó por debajo del mínimo. El más bajo tiene 7,3:1 y la norma pide 4,5:1. \\ \hline
Tamaños de letra & \cumple & El texto se puede ampliar al 200\,\% y al 400\,\% sin que se corte ni aparezca una barra para mover la página hacia los lados. Dos de cada tres textos miden 14 píxeles o más. \\ \hline
Navegación por teclado & \cumple & Todas las pantallas se pueden recorrer y usar solo con el teclado, y siempre se marca dónde está el teclado. \\ \hline
\end{tabular}
\end{center}

\newpage
\section{Qué se revisó y cómo}

\subsection*{Pantallas revisadas}

Se revisaron 9 pantallas del prototipo, en tamaño de escritorio (1440 píxeles de ancho) y en los dos temas de color.

\begin{center}
\begin{tabular}{|c|>{\raggedright\arraybackslash}m{6.2cm}|>{\raggedright\arraybackslash}m{8.2cm}|}
\hline
\rowcolor{barra}
\textbf{\#} & \textbf{Pantalla} & \textbf{Qué contiene} \\ \hline
1 & Iniciar sesión & Selector de perfil, correo y contraseña. \\ \hline
2 & Cortes de obra & Cifras, buscador y tabla de cortes. \\ \hline
3 & Detalle del corte & Datos del corte e historial. \\ \hline
4 & Registrar corte de obra & Formulario largo con fechas, actividades y personal. \\ \hline
5 & Personal del corte & Lista de trabajadores con casillas. \\ \hline
6 & Personal SST & Tarjetas de filtro y tabla de trabajadores. \\ \hline
7 & Ficha del trabajador & Cursos, aptitud médica y certificados. \\ \hline
8 & Registrar curso & Formulario con carga de certificado. \\ \hline
9 & Actualizar aptitud médica & Formulario con vista previa del semáforo. \\ \hline
\end{tabular}
\end{center}

\begin{center}
\begin{tabular}{@{}c@{\hspace{0.4cm}}c@{\hspace{0.4cm}}c@{}}
\escritorio{5.4cm}{acceso} & \escritorio{5.4cm}{cortes-claro} & \escritorio{5.4cm}{formulario} \\[2pt]
\pie{1. Iniciar sesión} & \pie{2. Cortes de obra} & \pie{4. Registrar corte} \\[8pt]
\escritorio{5.4cm}{detalle} & \escritorio{5.4cm}{sst-claro} & \escritorio{5.4cm}{ficha} \\[2pt]
\pie{3. Detalle del corte} & \pie{6. Personal SST} & \pie{7. Ficha del trabajador}
\end{tabular}
\end{center}

\subsection*{Cómo se hizo}

\begin{enumerate}
  \item \textbf{Medición del contraste.} Se tomó el color de cada texto visible y el color de su fondo, y el contraste se calculó en una página web con la fórmula de WCAG. Se hizo en tema claro y en tema oscuro.
  \item \textbf{Recorrido con el teclado.} Se recorrió cada pantalla solo con las teclas \tecla{Tab}, \tecla{Mayús}+\tecla{Tab}, \tecla{Enter}, \tecla{Espacio} y \tecla{Esc}, revisando el orden, que siempre se vea dónde está el teclado y que todo se pueda activar.
  \item \textbf{Ampliación del texto.} Se simuló el zoom del navegador al 200\,\% y al 400\,\% para ver si el contenido se sigue leyendo completo.
\end{enumerate}

\subsection*{Reglas que se usaron}

\begin{center}
\begin{tabular}{|>{\raggedright\arraybackslash}m{5.6cm}|>{\raggedright\arraybackslash}m{7.4cm}|>{\centering\arraybackslash}m{2.6cm}|}
\hline
\rowcolor{barra}
\textbf{Regla (WCAG 2.1)} & \textbf{Qué se pide} & \textbf{Mínimo} \\ \hline
1.4.3 Contraste del texto & Texto normal. & 4,5:1 \\ \hline
1.4.3 Contraste del texto & Texto grande: 24 píxeles o más, o 18,7 píxeles o más en negrita. & 3:1 \\ \hline
1.4.11 Contraste de controles & Bordes de los campos y recuadro que marca dónde está el teclado. & 3:1 \\ \hline
1.4.4 Cambio de tamaño del texto & Ampliar al 200\,\% sin perder contenido ni funciones. & 200\,\% \\ \hline
1.4.10 Ajuste del contenido & A 320 píxeles de ancho (zoom 400\,\%) no debe hacer falta mover la página hacia los lados. & 320 px \\ \hline
2.1.1 y 2.1.2 Teclado & Todo se puede usar con teclado y el teclado nunca queda atrapado. & No aplica \\ \hline
2.4.1, 2.4.3 y 2.4.7 Navegación & Atajo para saltar al contenido, orden lógico y foco siempre visible. & No aplica \\ \hline
\end{tabular}
\end{center}

\section{Contraste de color}

\textbf{Resultado: \cumple.} Ninguno de los 1.056 textos medidos quedó por debajo del mínimo. Todos los textos superan incluso 7:1, que es el nivel más exigente de la norma (AAA). Esto se debe a que los colores del prototipo se eligieron desde el sistema de diseño con buen contraste.

\subsection*{Tema claro}

\begin{center}
\begin{tabular}{|>{\raggedright\arraybackslash}m{5.3cm}|>{\centering\arraybackslash}m{3.5cm}|>{\centering\arraybackslash}m{1.8cm}|>{\centering\arraybackslash}m{1.5cm}|>{\raggedright\arraybackslash}m{2.2cm}|}
\hline
\rowcolor{barra}
\textbf{Dónde se usa} & \textbf{Muestra} & \textbf{Contraste} & \textbf{Mínimo} & \textbf{Resultado} \\ \hline
Texto principal sobre tarjetas & \muestra{0A1F2E}{FFFFFF}{Corte \#06} & 16,8:1 & 4,5:1 & \cumple \\ \hline
Texto principal sobre el fondo & \muestra{0A1F2E}{F4F8FB}{Iniciar sesión} & 15,8:1 & 4,5:1 & \cumple \\ \hline
Texto secundario (cédulas, fechas) & \muestra{1F4A66}{FFFFFF}{C.C. 16.789.432} & 9,4:1 & 4,5:1 & \cumple \\ \hline
Encabezados de las tablas & \muestra{1F4A66}{E8F0F6}{Corte / periodo} & 8,2:1 & 4,5:1 & \cumple \\ \hline
Botón principal & \muestra{FFFFFF}{0B3C5D}{Registrar corte} & 11,5:1 & 4,5:1 & \cumple \\ \hline
Botón secundario y enlaces & \muestra{0B3C5D}{FFFFFF}{Ver detalle} & 11,5:1 & 4,5:1 & \cumple \\ \hline
Menú lateral & \muestra{E8F2F8}{0B3C5D}{Componentes} & 10,2:1 & 4,5:1 & \cumple \\ \hline
Opción activa del menú \newline {\small(la más baja de todo el prototipo)} & \muestra{E8F2F8}{1A5279}{Cortes de obra} & 7,3:1 & 4,5:1 & \cumple \\ \hline
Estado ``Visado'' y ``Habilitado'' & \muestra{14532D}{ECFBF1}{Visado} & 8,5:1 & 4,5:1 & \cumple \\ \hline
Estado ``Devuelto'' y ``Vence pronto'' & \muestra{7A3A06}{FFF6E6}{Devuelto} & 8,0:1 & 4,5:1 & \cumple \\ \hline
Estado ``En revisión'' & \muestra{123F6E}{EAF3FC}{En revisión} & 9,5:1 & 4,5:1 & \cumple \\ \hline
Bloqueos y errores & \muestra{8B1A1A}{FFF0F0}{Bloqueado} & 8,4:1 & 4,5:1 & \cumple \\ \hline
\end{tabular}
\end{center}

\subsection*{Tema oscuro}

\begin{center}
\begin{tabular}{|>{\raggedright\arraybackslash}m{5.3cm}|>{\centering\arraybackslash}m{3.5cm}|>{\centering\arraybackslash}m{1.8cm}|>{\centering\arraybackslash}m{1.5cm}|>{\raggedright\arraybackslash}m{2.2cm}|}
\hline
\rowcolor{barra}
\textbf{Dónde se usa} & \textbf{Muestra} & \textbf{Contraste} & \textbf{Mínimo} & \textbf{Resultado} \\ \hline
Texto principal & \muestra{E8F2F8}{0C2331}{Corte \#06} & 14,2:1 & 4,5:1 & \cumple \\ \hline
Texto secundario & \muestra{A9CCE0}{0C2331}{C.C. 16.789.432} & 9,5:1 & 4,5:1 & \cumple \\ \hline
Botón principal & \muestra{061722}{8CCFF0}{Registrar corte} & 10,7:1 & 4,5:1 & \cumple \\ \hline
Opción activa del menú & \muestra{E8F2F8}{14344A}{Cortes de obra} & 11,4:1 & 4,5:1 & \cumple \\ \hline
Estado ``Visado'' & \muestra{6EE7A0}{0B2A1A}{Visado} & 10,0:1 & 4,5:1 & \cumple \\ \hline
Estado ``Devuelto'' & \muestra{FCC86B}{2E1F08}{Devuelto} & 10,3:1 & 4,5:1 & \cumple \\ \hline
Estado ``En revisión'' & \muestra{9CC9F5}{0D2740}{En revisión} & 8,7:1 & 4,5:1 & \cumple \\ \hline
Bloqueos & \muestra{FCA5A5}{331314}{Bloqueado} & 8,9:1 & 4,5:1 & \cumple \\ \hline
Mensajes de error \newline {\small(la más baja del tema oscuro)} & \muestra{FCA5A5}{0C2331}{No se pueden...} & 8,5:1 & 4,5:1 & \cumple \\ \hline
\end{tabular}
\end{center}

\begin{center}
\begin{tabular}{@{}c@{\hspace{0.6cm}}c@{}}
\escritorio{8.2cm}{cortes-claro} & \escritorio{8.2cm}{cortes-oscuro} \\[2pt]
\pie{Tema claro} & \pie{Tema oscuro}
\end{tabular}
\end{center}

\subsection*{Bordes de campos y recuadro de foco}

Además del texto, la norma pide que los bordes de los campos y el recuadro que muestra dónde está el teclado se distingan del fondo (mínimo 3:1).

\begin{center}
\begin{tabular}{|>{\raggedright\arraybackslash}m{5.3cm}|>{\centering\arraybackslash}m{3.5cm}|>{\centering\arraybackslash}m{1.8cm}|>{\centering\arraybackslash}m{1.5cm}|>{\raggedright\arraybackslash}m{2.2cm}|}
\hline
\rowcolor{barra}
\textbf{Elemento} & \textbf{Muestra} & \textbf{Contraste} & \textbf{Mínimo} & \textbf{Resultado} \\ \hline
Borde de campos sobre el fondo & \muestraborde{4F7A96}{F4F8FB} & 4,3:1 & 3:1 & \cumple \\ \hline
Borde de campos sobre tarjetas & \muestraborde{4F7A96}{FFFFFF} & 4,6:1 & 3:1 & \cumple \\ \hline
Borde de campos, tema oscuro & \muestraborde{6F9AB5}{0C2331} & 5,4:1 & 3:1 & \cumple \\ \hline
Recuadro de foco & \muestraborde{0B3C5D}{F4F8FB} & 10,8:1 & 3:1 & \cumple \\ \hline
Recuadro de foco en el menú & \muestraborde{8CCFF0}{0B3C5D} & 6,8:1 & 3:1 & \cumple \\ \hline
Recuadro de foco, tema oscuro & \muestraborde{8CCFF0}{061722} & 10,7:1 & 3:1 & \cumple \\ \hline
\end{tabular}
\end{center}

\newpage
\section{Tamaños de letra}

\textbf{Resultado: \cumple.} La norma no fija un tamaño mínimo de letra, pero sí pide que el texto se pueda ampliar al 200\,\% sin perder contenido. El prototipo cumple y, además, usa tamaños cómodos de leer.

\subsection*{Tamaños usados}

Conteo de los 528 textos medidos en las 9 pantallas (tema claro).

\begin{center}
\begin{tabular}{|>{\centering\arraybackslash}m{1.8cm}|>{\raggedright\arraybackslash}m{9.6cm}|>{\centering\arraybackslash}m{1.5cm}|>{\centering\arraybackslash}m{1.6cm}|}
\hline
\rowcolor{barra}
\textbf{Tamaño} & \textbf{Dónde se usa} & \textbf{Textos} & \textbf{\%} \\ \hline
36 px & Frase de bienvenida en ``Iniciar sesión''. & 1 & 0,2\,\% \\ \hline
28 px & Título de cada pantalla y cifras grandes de las tarjetas. & 16 & 3,0\,\% \\ \hline
18 px & Títulos de las secciones de un formulario o ficha. & 15 & 2,8\,\% \\ \hline
16 px & Texto dentro de los campos, descripciones y barra superior. & 52 & 9,8\,\% \\ \hline
\rowcolor{suave}
14 px & \textbf{Texto general:} etiquetas, botones, menú y celdas principales. & 270 & 51,1\,\% \\ \hline
12 px & Encabezados de tabla, fechas, estados y datos secundarios. & 173 & 32,8\,\% \\ \hline
11 px & Rótulo ``OPERACIONES Y CALIDAD'' en ``Iniciar sesión'' (decorativo). & 1 & 0,2\,\% \\ \hline
\end{tabular}
\end{center}

\begin{itemize}
  \item El texto de los \textbf{campos} mide 16 píxeles. En el celular este tamaño evita que el teléfono haga zoom solo al tocar un campo.
  \item Los títulos tienen un tamaño claramente mayor que el texto, lo que ayuda a ubicarse en la pantalla.
\end{itemize}

\subsection*{Prueba de ampliación (zoom)}

Se simuló una ventana de 1280 píxeles con zoom al 200\,\% y al 400\,\%.

\begin{center}
\begin{tabular}{@{}m{4.9cm}@{\hspace{0.45cm}}m{4.9cm}@{\hspace{0.5cm}}m{6.6cm}@{}}
\escritorio{4.9cm}{zoom-200} & \escritorio{4.9cm}{zoom-200-formulario} &
\begin{itemize}
  \item \textbf{Todo se lee completo:} nada se corta ni se monta encima de otra cosa.
  \item \textbf{La página no se sale de la pantalla:} no aparece barra para moverse hacia los lados, ni al 200\,\% ni al 400\,\%.
  \item El menú se guarda en el botón de tres líneas, igual que en el celular.
  \item La tabla de cortes se desplaza \emph{dentro de su propio recuadro}. Esto está permitido para tablas.
\end{itemize} \\[2pt]
\centering\pie{Cortes de obra al 200\,\%} & \centering\pie{Registrar corte al 200\,\%} &
\end{tabular}
\end{center}

\newpage
\section{Navegación por teclado}

\textbf{Resultado: \cumple.} Todas las pantallas se pueden usar sin ratón.

\recuadro{Cómo se usa el teclado en una página web}{%
\tecla{Tab} pasa al siguiente elemento - \tecla{Mayús}+\tecla{Tab} vuelve al anterior - \tecla{Enter} o \tecla{Espacio} activa botones y casillas - \tecla{Esc} cierra paneles y menús. El \textbf{foco} es el elemento donde está ``parado'' el teclado; en Loard se marca con un recuadro azul alrededor.}

\subsection*{Lista de revisión}

\begin{center}
\begin{tabular}{|c|>{\raggedright\arraybackslash}m{12.2cm}|>{\raggedright\arraybackslash}m{2.4cm}|}
\hline
\rowcolor{barra}
\textbf{\#} & \textbf{Qué se revisó} & \textbf{Resultado} \\ \hline
1 & Con el primer \tecla{Tab} aparece el botón ``Saltar al contenido''. Con \tecla{Enter} se llega directo al contenido, sin pasar por todo el menú. & \cumple \\ \hline
2 & El orden sigue la lectura: menú, barra superior y luego el contenido, de izquierda a derecha y de arriba abajo. & \cumple \\ \hline
3 & Siempre se ve dónde está el teclado. & \cumple \\ \hline
4 & Botones, enlaces, listas desplegables, casillas y tarjetas de filtro se activan con \tecla{Enter} o \tecla{Espacio}. & \cumple \\ \hline
5 & Panel del certificado y menú del celular: mientras están abiertos el teclado se queda dentro, \tecla{Esc} los cierra y el foco vuelve al botón que los abrió. & \cumple \\ \hline
6 & No hay trampas: después del último elemento, \tecla{Tab} sale de la página y vuelve a empezar. & \cumple \\ \hline
7 & Una tabla que no cabe se puede mover con las flechas, porque su recuadro también recibe el foco. & \cumple \\ \hline
8 & Los trabajadores bloqueados se alcanzan con \tecla{Tab} y se lee el motivo, pero su casilla no se marca con \tecla{Espacio}. & \cumple \\ \hline
9 & Al guardar un corte incompleto, el aviso aparece arriba del formulario y la página sube para mostrarlo. & \cumple \\ \hline
\end{tabular}
\end{center}

\newpage
\subsection*{Así se ve el foco}

\begin{center}
\begin{tabular}{@{}>{\centering\arraybackslash}m{8.1cm}@{\hspace{0.6cm}}>{\centering\arraybackslash}m{8.1cm}@{}}
\recorte{7.2cm}{teclado-saltar} & \recorte{7.2cm}{teclado-filtro} \\[1pt]
\pie{``Saltar al contenido'', primer Tab} & \pie{Tarjeta de filtro en Personal SST} \\[6pt]
\recorte{7.2cm}{teclado-menu} & \recorte{7.2cm}{teclado-boton} \\[1pt]
\pie{Opción del menú} & \pie{Botón principal} \\[6pt]
\recorte{7.2cm}{teclado-campo} & \recorte{7.2cm}{teclado-oscuro} \\[1pt]
\pie{Campo de búsqueda} & \pie{Botón principal en tema oscuro} \\[6pt]
\recorte{7.2cm}{teclado-casilla} & \recorte{7.2cm}{teclado-bloqueado} \\[1pt]
\pie{Casilla marcada con Espacio} & \pie{Trabajador bloqueado: no se marca} \\
\end{tabular}

\smallskip
\recorte{15.6cm}{teclado-tabla}\\[1pt]
\pie{Botón ``Ver detalle'' dentro de la tabla}
\end{center}

\subsection*{Paneles y menú del celular}

\begin{center}
\begin{tabular}{@{}>{\centering\arraybackslash}m{10.4cm}@{\hspace{0.6cm}}>{\centering\arraybackslash}m{4.4cm}@{}}
\escritorio{9.2cm}{teclado-panel} & \movil{6.0cm}{teclado-menu-movil} \\[2pt]
\pie{Panel del certificado abierto con Enter} & \pie{Menú del celular abierto con Enter}
\end{tabular}
\end{center}

En ambos, \tecla{Tab} da vueltas solo dentro del panel, \tecla{Esc} lo cierra y el foco vuelve al botón que lo abrió.

\end{document}
